
AI alignment methods typically optimize models through Reinforcement Learning from Human Feedback (RLHF), treating alignment as unidirectional adjustment of AI policy to human preferences. This framing does not capture that conversations are bidirectional interactions in which users also adapt their communication strategies. This paper proposes behavioral coupling---quantified as correlation between user learning rate (reformulation slope) and AI responsiveness (query-response diversity correlation)---as a metric for bidirectional alignment. Analysis of 169,352 multi-turn conversations from the HH-RLHF dataset demonstrates that users reformulate queries less frequently over turns (mean slope = -0.021, p=0.012, Cohen's d=-0.246) and AI response diversity correlates with query diversity (r=0.396, p<0.001, 95\% CI [0.392, 0.400]). Correlation strengthens with conversation length (r=0.04 for 2-3 turns to r=0.19 for 6+ turns), consistent with alignment emerging through extended interaction. Effect sizes are smaller than initially hypothesized and the core coupling hypothesis remains untested. These findings validate the theoretical foundation of bidirectional alignment and demonstrate that co-adaptation signals can be extracted from conversation logs without additional human evaluation.
